\subsection{Research Methodology}

This project uses a hands-on research and development approach. We combined theoretical research with system design and practical software implementation to build the platform.

The work was divided into the following key phases:

\begin{itemize}

    \item \textbf{Literature Review:} We started by studying academic papers, technical journals, and documentation on blockchain, smart contracts, and cryptographic hashing. This helped us understand the limitations of current setups and identify new technological options.

    \item \textbf{Requirement Analysis:} Next, we identified the functional and non-functional requirements of the system based on the real-world flaws found in traditional verification methods. This included analyzing user roles, security needs, and backend workflows to shape our system architecture.

    \item \textbf{System Design:} We mapped out the entire architecture for the NFT platform. This phase involved structuring the blockchain integration, planning the database schema, designing user interfaces, and mapping out the step-by-step verification flows.

    \item \textbf{Smart Contract Development:} We wrote and deployed the actual smart contracts responsible for handling the core logic of the system, including minting the certificate NFTs, assigning ownership, and managing the verification rules on-chain.

    \item \textbf{Hashing Implementation:} We integrated the SHA-256 cryptographic hashing algorithm into the backend. This allows the system to generate a unique digital footprint for every certificate file, which acts as the main line of defense against document tampering.

    \item \textbf{System Validation:} This phase focused on verifying the core functions of the deployed smart contracts and the web interface. Validation involved running sample certificate files through the SHA-256 hashing pipeline and ensuring that the generated hashes successfully matched the records on the blockchain ledger.

    \item \textbf{Deployment Review:} The final stage involved analyzing the system's behavior during execution. This review confirmed that the platform successfully automates certificate verification and flags modified documents instantly, demonstrating a practical alternative to slow, centralized methods.
\end{itemize}